% Part: methods
% Chapter: proofs
% Section: using-definitions

\documentclass[../../../include/open-logic-section]{subfiles}

\begin{document}

% BN-SRC-476: the part identifier is mth, not mod.
\olfileid{mth}{prf}{def}

\olsection{সংজ্ঞার ব্যবহার}

বলেছি, প্রমাণে লাগতে পারে এমন সব সংজ্ঞা জানতে এবং সেগুলি
সঠিকভাবে প্রয়োগ করতে পারতে হবে। বিষয়টি এত গুরুত্বপূর্ণ যে
আরও একটু খুঁটিয়ে দেখা দরকার। কোনো ধর্ম বা সম্পর্ককে সংক্ষেপে
উল্লেখ করার জন্য সংজ্ঞা দিই। যে সংক্ষিপ্ত পদটি চালু করি,
তাকে \emph{সংজ্ঞেয় পদ} (definiendum) এবং পদটি দিয়ে যা বোঝাই,
তাকে \emph{সংজ্ঞাদায়ী অংশ} (definiens) বলি। প্রমাণে প্রায়ই
সংজ্ঞেয় পদ কীভাবে চালু হয়েছে, সেখানে ফিরে যেতে হয়। কারণ
প্রমাণ এগিয়ে নিতে সংজ্ঞাদায়ী অংশের—অর্থাৎ সংক্ষিপ্ত পদটি
যার বদলে বসেছে সেই বিস্তারিত কথার—যৌক্তিক গঠন কাজে
লাগাতে হয়। সংজ্ঞা খুলে লিখলে যুক্তির আসল কাজটি কোথায়,
তা বোঝা যায়।

একটি উদাহরণ দিয়ে শুরু করি। ধরো, প্রমাণ করতে চাও:

\begin{prop}
যেকোনো সেট $A$ ও $B$-র জন্য $A \cup B = B \cup A$।
\end{prop}

প্রমাণ শুরু করতেই জানতে হবে, দুটি সেট অভিন্ন হওয়ার অর্থ
কী; অর্থাৎ সেটের ক্ষেত্রে ওই সমীকরণের ``$=$'' কী বোঝায়।
সেটদুটির !!{element}s একই হলে এবং কেবল তখনই তাদের অভিন্ন
বলি। তাই যে সংজ্ঞাটি খুলে লিখতে হবে, তা হলো:

\begin{defn}
সেট $A$ ও $B$ \emph{অভিন্ন}, $A = B$, তখনই, যখন
$A$-র প্রত্যেক !!{element} $B$-রও !!a{element}
এবং উল্টোটিও সত্য।
\end{defn}

এখানে $A$ ও $B$ যেকোনো সেটের স্থানধারক। সংজ্ঞাটি
যাকে সংজ্ঞা দেয়—\emph{সংজ্ঞেয় পদ}—তা হলো
``$A = B$'' রাশিটি; $A = B$ সত্য হওয়ার শর্ত বলেই
সংজ্ঞাটি কাজ করে। সেই শর্ত—``$A$-র প্রত্যেক
!!{element} $B$-রও !!a{element}, এবং উল্টোটিও''—
হলো \emph{সংজ্ঞাদায়ী অংশ}।\footnote{এই বিশেষ
ক্ষেত্রে—বিভ্রান্তি এখানেই!—$A = B$ হলে $A$ ও $B$
আসলে একই সেট, যদিও সমীকরণের দুই পাশে তার জন্য
দুটি আলাদা অক্ষর লিখি। কিন্তু একই সেটকে শনাক্ত করার
পদ্ধতি আলাদা হতে পারে; তাই সংজ্ঞাটি তুচ্ছ নয়।}
সংজ্ঞাটি বলে, শর্তটি সত্য হলে \emph{এবং কেবল তখনই}
$A = B$ সত্য (ইংরেজিতে ``iff'' সংক্ষিপ্ত রূপটি
ব্যবহার করা হয়)।

সংজ্ঞা প্রয়োগের সময় তার $A$ ও $B$-কে আলোচ্য
ক্ষেত্রের সঙ্গে মিলাতে হবে। এখানে $A \cup B = B \cup A$
সত্য দেখাতে হলে প্রত্যেক $z \in A \cup B$-কে
$B \cup A$-তেও থাকতে হবে, এবং উল্টোটিও দেখাতে
হবে। প্রস্তাবের $A \cup B$ সংজ্ঞার $A$-র ভূমিকা,
আর প্রস্তাবের $B \cup A$ সংজ্ঞার $B$-র ভূমিকা পালন
করে। একই অক্ষর $A$, $B$ সংজ্ঞা ও প্রমাণাধীন
প্রস্তাবে ভিন্ন ভূমিকায় এসেছে—এ দুটি গুলিয়ে ফেলো না।
যেমন, $A$-র প্রত্যেক !!{element} $B$-রও
!!a{element}, এবং উল্টোটিও দেখালেই প্রস্তাবটি
প্রমাণিত হবে ভাবা ভুল। তাতে $A = B$ প্রমাণ হবে,
$A \cup B = B \cup A$ নয়। (তাছাড়া $A$ ও $B$
যেকোনো দুটি সেট হতে পারে; তাদের সম্বন্ধে আর কিছু
না ধরলে $A$ ও $B$ ভিন্নও হতে পারে।)

প্রমাণে সংযোগের মতো সেটতাত্ত্বিক ধারণা আছে। তাই
প্রমাণ কীভাবে এগোবে বুঝতে $\cup$-র অর্থও জানতে হবে।
কখনো একটি সংজ্ঞা খুললে আরও সংজ্ঞা খোলার দরকার হয়।
যেমন, $A \cup B$-র সংজ্ঞা হলো
$\Setabs{z}{z \in A \text{ অথবা } z \in B}$।
তাই $x \in A \cup B$ দেখাতে চাইলে $\cup$-র
সংজ্ঞা খুলে দেখতে হয়,
$x \in \Setabs{z}{z \in A \text{ অথবা } z \in B}$
দেখাতে হবে। এবার মনে রাখতে হবে,
$x \in \Setabs{z}{\dots z\dots}$ তখনই,
যখন $\dots x\dots$। সুতরাং
$\Setabs{z}{\dots z \dots}$ সংকেতটিও খুললে
যা দেখাতে হয় তা হলো: $x \in A$ অথবা $x \in B$।
তাই ``$A \cup B$-র প্রত্যেক !!{element}
$B \cup A$-রও !!a{element}'' কথাটির পূর্ণ অর্থ:
``প্রত্যেক $x$-এর জন্য, $x \in A$ অথবা $x \in B$
হলে $x \in B$ অথবা $x \in A$।'' প্রস্তাবের সব
সংজ্ঞা খুললে দেখতে পাই, আসলে দেখাতে হবে:

\begin{prop}
যেকোনো সেট $A$ ও $B$-র জন্য: (a) প্রত্যেক $x$-এর
জন্য, $x \in A$ অথবা $x \in B$ হলে $x \in B$
অথবা $x \in A$; এবং (b) প্রত্যেক $x$-এর জন্য,
$x \in B$ অথবা $x \in A$ হলে $x \in A$
অথবা $x \in B$।
\end{prop}

প্রমাণ গড়তে সংজ্ঞা খুলে দেখা অপরিহার্য—এটিই
মূল কথা। ঠিকভাবে তা করা কখনো কঠিন: সংজ্ঞার
চলক ও প্রমাণাধীন দাবির পদগুলিকে আলাদা করে
চিনতে এবং যথাযথভাবে মিলাতে হয়। সফল হতে
প্রশ্নটি কী চায় এবং তার সব পদের অর্থ কী,
তা বুঝতে হবে—প্রায়ই একাধিক সংজ্ঞা খুলতে হয়।
নিচের মতো সরল প্রমাণে সংজ্ঞা থেকেই প্রায়
সঙ্গে সঙ্গে সমাধান মেলে। অবশ্য সব সময় এত
সহজ হবে না।

\begin{prob}
ধরো, তোমাকে $A \cap B \neq \emptyset$ প্রমাণ করতে
বলা হলো। এখানে ব্যবহৃত সব সংজ্ঞা খুলে কথাটি
আবার বলো; অর্থাৎ এমনভাবে লেখো যাতে
``$\cap$'', ``='', বা ``$\emptyset$'' না আসে।
\end{prob}

\end{document}
